\chapter{Supplementary Mathematical Results}
\label{app:chow-iii-selected-exercise-solutions}
\dcref{appJ}
\source{chow-et-al-ricci-flow-part-iii:page-0518,chow-et-al-ricci-flow-part-iii:page-0519,chow-et-al-ricci-flow-part-iii:page-0520,chow-et-al-ricci-flow-part-iii:page-0521,chow-et-al-ricci-flow-part-iii:page-0522}

\section{Ricci-flow results}

\begin{proposition}[Point-picking doubling lemma]
\label{prop:chow-iii-appj-solution-22-2}
\source{chow-et-al-ricci-flow-part-iii:page-0205,chow-et-al-ricci-flow-part-iii:page-0206,chow-et-al-ricci-flow-part-iii:page-0518}
\uses{def:chow-iii-alpha-small-large-curvature}
\dcref{appJ:22.2}
\group{Point picking}


Let \((\mathcal M^n,g(t))\), \(t\in[0,T]\), have bounded curvature, let
\(\alpha>0\), \(0<\epsilon<1\), and suppose
\((x_0,t_0)\in\mathcal M\times[T/2,T]\) satisfies
\[
 Q=|\Rm|(x_0,t_0)>\max\{\alpha/t_0,8/(T\epsilon)\}.
\]
Then the doubling point-picking procedure terminates at an
\(\alpha\)-large point \((\widehat x_0,\widehat t_0)\), where
\(\widehat t_0\geq T/4\), \(\widehat Q=|\Rm|(\widehat x_0,\widehat t_0)\geq Q\),
and
\[
 |\Rm|(x,t)\leq2\widehat Q
\]
throughout
\(P(\widehat x_0,\widehat t_0,
(\epsilon\widehat Q)^{-1/2},-(\epsilon\widehat Q)^{-1})\).
\end{proposition}
\begin{proof}
\sketch
Set \(H=TQ/8\).  If the current point \((x_{j-1},t_{j-1})\) does not
have the required curvature control, choose a point in its backward
parabolic cylinder for which
\[
 Q_j=|\Rm|(x_j,t_j)>2Q_{j-1}.
\]
Thus \(Q_j>2^jQ\).  Since \(H>\epsilon^{-1}\), the time coordinates satisfy
\[
 t_j\geq t_{j-1}-HQ_{j-1}^{-1}
 \geq t_0-HQ^{-1}\sum_{k=0}^{j-1}2^{-k}
 \geq \frac T4(1+2^{-j}).
\]
In particular, every point remains in \(\mathcal M\times[T/4,T]\).  This
estimate also gives \(t_j\geq t_{j-1}/2\).  Inductively, therefore,
\[
 Q_j>2Q_{j-1}>\frac{2\alpha}{t_{j-1}}
 \geq\frac{\alpha}{t_j},
\]
so every selected point remains \(\alpha\)-large.  An infinite sequence would
have \(Q_j\to\infty\), contrary to the assumed uniform curvature bound on
\(\mathcal M\times[0,T]\).  Hence the procedure stops, and the stopping
condition is precisely the asserted local curvature estimate.
\end{proof}

\begin{proposition}[Point-picking selection lemma]
\label{prop:chow-iii-appj-solution-22-5}
\source{chow-et-al-ricci-flow-part-iii:page-0209,chow-et-al-ricci-flow-part-iii:page-0518,chow-et-al-ricci-flow-part-iii:page-0519}
\uses{lem:chow-iii-point-picking-subset}
\dcref{appJ:22.5}
\group{Point picking}


Let \((\mathcal M^n,g(t))\), \(t\in[0,\epsilon^2]\), be a continuous family
of complete smooth metrics, let \(\mathcal S\subset\mathcal M\times
(0,\epsilon^2]\), and suppose \((x_1,t_1)\in\mathcal S\) satisfies
\[
 d_{g(t_1)}(x_1,x_0)<\epsilon_0,
 \qquad |\Rm|(x_1,t_1)>\epsilon_0^{-2}.
\]
For every \(A>0\) there is \((\bar x,\bar t)\in\mathcal S\), with
\(0<\bar t\leq t_1\) and
\(d_{g(\bar t)}(\bar x,x_0)<(2A+1)\epsilon_0\), such that
\(|\Rm|(\bar x,\bar t)>\epsilon_0^{-2}\) and
\[
 |\Rm|(x,t)\leq4|\Rm|(\bar x,\bar t)
\]
whenever \((x,t)\in\mathcal S\), \(0<t\leq\bar t\), and
\[
 d_{g(t)}(x,x_0)\leq d_{g(\bar t)}(\bar x,x_0)
     +A|\Rm|(\bar x,\bar t)^{-1/2}.
\]
\end{proposition}
\begin{proof}
\sketch
Starting with \((x_1,t_1)\), if the desired conclusion fails at
\((x_m,t_m)\), choose \((x_{m+1},t_{m+1})\in\mathcal S\) such that
\[
 |\Rm|(x_{m+1},t_{m+1})>4|\Rm|(x_m,t_m),
 \qquad 0<t_{m+1}\leq t_m,
\]
and
\[
 d_{g(t_{m+1})}(x_{m+1},x_0)
 \leq d_{g(t_m)}(x_m,x_0)+A|\Rm|(x_m,t_m)^{-1/2}.
\]
Writing \(Q_k=|\Rm|(x_k,t_k)\), induction gives
\(Q_k>4^{k-1}Q_1\), and hence
\(Q_k^{-1/2}<2^{1-k}Q_1^{-1/2}\).  Consequently,
\[
 d_{g(t_k)}(x_k,x_0)
 <\epsilon_0+A\sum_{j=1}^{k-1}Q_j^{-1/2}
 <\epsilon_0+2AQ_1^{-1/2}
 <(2A+1)\epsilon_0.
\]
Thus all selected points lie in the indicated space-time tube.  If the
construction never stopped, \(Q_k\to\infty\).  Completeness and continuity
of the metrics on the compact time interval make the closed bounded tube
compact, so \(|\Rm|\) is uniformly bounded there, a contradiction.  At the
first stopping index the stated factor-four bound is exactly the stopping
condition.
\end{proof}

\begin{proposition}[Symmetrization layer-cake identity]
\label{prop:chow-iii-appj-solution-22-18}
\source{chow-et-al-ricci-flow-part-iii:page-0230,chow-et-al-ricci-flow-part-iii:page-0232,chow-et-al-ricci-flow-part-iii:page-0519}
\uses{thm:chow-iii-log-sobolev-isoperimetric}
\dcref{appJ:22.18}
\group{Logarithmic Sobolev}


Let \(h:\mathbb R^n\to[0,\infty)\) be the radial rearrangement used in
spherical symmetrization, let
\(F(t)=\Vol\{y:h(y)\geq t\}\), and let \(\lambda\) be absolutely continuous
with vanishing boundary terms.  Then
\[
 \int_0^\infty\lambda'(s)F(s)\,ds
 =\int_{\{h>0\}}\lambda(h)\,d\mu_{\mathbb R^n}.
\]
\end{proposition}
\begin{proof}
\sketch
The co-area formula gives
\[
 F(t)=\int_t^\infty\left(\int_{\{h=s\}}|\nabla h|^{-1}\,d\sigma\right)ds,
 \qquad
 -F'(s)=\int_{\{h=s\}}|\nabla h|^{-1}\,d\sigma
\]
for almost every \(s>0\).  Integration by parts, followed by the co-area
formula once more, yields
\[
 \int_0^\infty\lambda'(s)F(s)\,ds
 =-\int_0^\infty\lambda(s)F'(s)\,ds
 =\int_0^\infty\lambda(s)
   \int_{\{h=s\}}|\nabla h|^{-1}\,d\sigma\,ds
 =\int_{\{h>0\}}\lambda(h)\,d\mu_{\mathbb R^n}.
\]
\end{proof}

\begin{proposition}[Nested-cylinder Moser estimate]
\label{prop:chow-iii-appj-solution-25-7}
\source{chow-et-al-ricci-flow-part-iii:page-0335,chow-et-al-ricci-flow-part-iii:page-0519,chow-et-al-ricci-flow-part-iii:page-0520}
\uses{ex:chow-iii-rescaled-moser-cylinder-iteration}
\dcref{appJ:25.7}
\group{Moser iteration}


For \(r\in[r_0,2r_0)\) and \(1<\gamma\leq2\), with
\(\gamma r\leq2r_0\), the nested-cylinder Moser iteration gives
\[
 \|v\|_{L^\infty(P_r)}\leq\widehat C(\gamma)
 \|v\|_{L^2(P_{\gamma r})},
\]
where
\[
 \widehat C(\gamma)=
 \frac{e^{C_S(1+\gamma\sqrt K r)n/4}}
 {r\Vol_{\widetilde g}(B_{\widetilde g}(x_0,r))^{1/2}}
 \widetilde C_0^{(n+3)n/4}\widehat C_n(\gamma-1)^{-n/4},
\]
and \(\widehat C_n\) depends only on \(n\).
\end{proposition}
\begin{proof}
\sketch
Set
\[
 r_i=r\left(1+\frac{\gamma-1}{2^{i-1}}\right).
\]
Then
\[
 r_i-r_{i+1}=\frac{\gamma-1}{2^i}r,
 \qquad
 r_i^2-r_{i+1}^2\geq\frac{\gamma-1}{2^{i-1}}r^2.
\]
The corresponding temporal and spatial cutoffs may be chosen so that
\[
 0\leq\partial_\tau\psi_i\leq
 \frac{2^i}{(\gamma-1)r^2},
 \qquad
 |\nabla\psi_i|_{\widetilde g}\leq
 \frac{2^{i+1}}{(\gamma-1)r},
\]
and therefore
\[
 \psi_i\partial_\tau\psi_i+|\nabla\psi_i|_{\widetilde g}^2
 \leq\frac{4^{i+2}}{(\gamma-1)r^2}.
\]
Insert these estimates into the one-step Sobolev inequality with
\(p_i=((n+2)/n)^{i-1}\).  This gives
\[
 \|v\|_{L^{2p_{i+1}}(P_{r_{i+1}})}
 \leq M_i^{1/(2p_{i+1})}
 \|v\|_{L^{2p_i}(P_{r_i})},
\]
where the factors \(M_i\) differ from those in the \(\gamma=2\) iteration
only by replacing the cutoff contribution with
\(4^{i+2}((\gamma-1)r^2)^{-1}\) and replacing \(r_i\) by the displayed
radii.  Multiplying over \(i\) is legitimate because, for
\(q=(n+2)/n\), both \(\sum q^{-i}\) and \(\sum iq^{-i}\) converge.
The identities \(\sum_{i\geq1}q^{-i}=n/2\) and the ball-volume comparison
produce the factor \((\gamma-1)^{-n/4}\), the power
\(\widetilde C_0^{(n+3)n/4}\), and the stated exponential and volume
factors.  Letting \(i\to\infty\) proves the estimate.
\end{proof}

\section{Metric-geometric results}

\begin{proposition}[Pullback pseudometric]
\label{prop:chow-iii-appj-solution-g-3}
\source{chow-et-al-ricci-flow-part-iii:page-0411,chow-et-al-ricci-flow-part-iii:page-0520}
\uses{prop:chow-iii-g3-pullback-metric}
\dcref{appJ:G.3}
\group{Metric spaces}


For a map \(f:X\to(Y,d_Y)\), the function
\((f^*d_Y)(x_1,x_2)=d_Y(f(x_1),f(x_2))\) is a pseudometric on \(X\), and it
is a metric if \(f\) is injective.
\end{proposition}
\begin{proof}
\sketch
The function is nonnegative and symmetric because \(d_Y\) is.  For any
three points \(x_1,x_2,x_3\in X\),
\[
 (f^*d_Y)(x_1,x_3)
 \leq(f^*d_Y)(x_1,x_2)+(f^*d_Y)(x_2,x_3)
\]
by the triangle inequality in \(Y\).  Finally, if \(f\) is injective and
\(x_1\neq x_2\), then \(f(x_1)\neq f(x_2)\), so
\((f^*d_Y)(x_1,x_2)>0\).
\end{proof}

\begin{example}[Nonunique tangent-cone limits]
\label{ex:chow-iii-appj-solution-g-24}
\source{chow-et-al-ricci-flow-part-iii:page-0419,chow-et-al-ricci-flow-part-iii:page-0520,chow-et-al-ricci-flow-part-iii:page-0521}
\uses{ex:chow-iii-g24-g28-nonunique-cones}
\dcref{appJ:G.24}
\group{Metric cones}


There is a connected compact intrinsic metric space \((X,d)\) with a point
\(p\) and two sequences \(\alpha_i,\beta_i\to\infty\) such that
\[
 (X,\alpha_i d,p)\longrightarrow\mathbb R^2,
 \qquad
 (X,\beta_i d,p)\longrightarrow[0,\infty)
\]
in pointed Gromov--Hausdorff distance.
\end{example}
\begin{proof}
\sketch
In the \((r,z)\)-half-plane let
\[
 \begin{split}
 Y={}&\{(0,0)\}
 \cup\bigcup_{k\geq0}\left([1/(2k+1)!,1/(2k)!]\times\{0\}\right)\\
 &\cup\left(\{0\}\times
 \bigcup_{k\geq0}[1/(2k+2)!,1/(2k+1)!]\right)\\
 &\cup\bigcup_{m\geq0}
 \{(t/m!,(1-t)/m!):0\leq t\leq1\}.
 \end{split}
\]
Revolve \(Y\) about the \(z\)-axis and give the resulting subset \(X\) of
\(\mathbb R^3\) its intrinsic metric.  The diagonal segments connect the
successive horizontal annuli and vertical axis intervals, so \(X\) is
connected.  It is closed and bounded, and its intrinsic diameter is finite
because the lengths of the successive pieces are bounded by a constant
multiple of \(\sum_{m\geq0}1/m!\); hence \(X\) is compact.

Take \(p=(0,0,0)\) and
\[
 \alpha_i=\sqrt{(2i+1)!(2i)!},
 \qquad
 \beta_i=\sqrt{(2i)!(2i-1)!}.
\]
After scaling by \(\alpha_i\), the horizontal annulus with radii
\(1/(2i+1)!\) and \(1/(2i)!\) has inner radius
\(1/\sqrt{2i+1}\) and outer radius \(\sqrt{2i+1}\).  Thus it exhausts the
plane on bounded sets, while every other feature is pushed either to the
basepoint or beyond each bounded ball.  The rescaled intrinsic metrics
therefore converge on bounded balls to the Euclidean plane.  Similarly, after
scaling by \(\beta_i\), the vertical interval with endpoints
\(1/(2i)!\) and \(1/(2i-1)!\) has scaled endpoints
\(1/\sqrt{2i}\) and \(\sqrt{2i}\); it exhausts a half-line and the remaining
features disappear from bounded pointed balls.  These bounded-ball
approximations give the two asserted pointed Gromov--Hausdorff limits, which
are not isometric.
\end{proof}

\begin{example}[Normalized gradient curve on a square]
\label{ex:chow-iii-appj-solution-h-44}
\source{chow-et-al-ricci-flow-part-iii:page-0465,chow-et-al-ricci-flow-part-iii:page-0521}
\uses{ex:chow-iii-h44-square-integral-curves}
\dcref{appJ:H.44}
\group{Integral curves}


Let \(C=[-1,1]^2\) and \(f(x,y)=d((x,y),\partial C)\).  The maximal integral
curve of \(\nabla f/|\nabla f|^2\) from \((x,y)\) first runs parallel to a
coordinate axis to the nearer diagonal \(y=x\) or \(y=-x\), and then follows
that diagonal to the origin.
\end{example}
\begin{proof}
\sketch
Since \(f(x,y)=1-\max\{|x|,|y|\}\), in each region
\(|x|>|y|\) or \(|y|>|x|\) the gradient points parallel to the coordinate
axis which decreases the larger absolute coordinate.  The normalized-gradient
curve consequently reaches the closest of the two diagonals without changing
the smaller absolute coordinate.  On a diagonal the generalized gradient is
the equal-weight combination of the two adjacent coordinate gradients; its
normalized integral curve stays on the diagonal and moves toward the origin.
The origin is the maximum set of \(f\), so this curve is maximal.
\end{proof}

\begin{example}[Unbounded Busemann function]
\label{ex:chow-iii-appj-solution-i-14}
\source{chow-et-al-ricci-flow-part-iii:page-0482,chow-et-al-ricci-flow-part-iii:page-0521,chow-et-al-ricci-flow-part-iii:page-0522}
\uses{def:chow-iii-i1-rays-busemann, def:chow-iii-g29-euclidean-cone}
\dcref{appJ:I.14}
\group{Busemann geometry}


There is a complete noncompact Riemannian surface of positive sectional
curvature carrying a ray whose Busemann function is unbounded below.
\end{example}
\begin{proof}
\sketch
Fix \(\alpha\in(\pi,2\pi)\) and let \(S^1_\alpha\) be the circle of length
\(\alpha\).  Choose opposite points \(\theta_1,\theta_2\), so
\(d(\theta_1,\theta_2)=\alpha/2\), and put
\(\beta=\cos(\alpha/2)<0\).  On the flat cone
\(\operatorname{Cone}(S^1_\alpha)\), for the ray
\(\gamma(s)=[(\theta_2,s)]\), the cone distance formula gives
\[
 \begin{split}
 b_\gamma([\theta_1,r])
 &=\lim_{s\to\infty}
 \left(s-\sqrt{r^2+s^2-2\beta rs}\right)\\
 &=\beta r,
 \end{split}
\]
which tends to \(-\infty\) as \(r\to\infty\).

Now take a complete positively curved rotationally symmetric smooth surface
asymptotic to this cone, for example the expanding gradient surface with cone
angle \(\alpha\).  Its rescaled distance functions converge on compact annuli
to those of the cone.  Hence, for the rays corresponding to \(\theta_1\) and
\(\theta_2\),
\[
 \frac{b_\gamma(\theta_1,r)}r\longrightarrow\beta<0
 \qquad(r\to\infty).
\]
Indeed, this follows by first evaluating the defining difference
\(s-d((\theta_1,r),\gamma(s))\) with \(s/r\) fixed, passing to the asymptotic
cone, and then letting \(s/r\to\infty\).  Thus the Busemann function on the
smooth complete positively curved surface is unbounded below.
\end{proof}

\begin{proposition}[Toponogov radial estimate]
\label{prop:chow-iii-appj-solution-i-38}
\source{chow-et-al-ricci-flow-part-iii:page-0498,chow-et-al-ricci-flow-part-iii:page-0500,chow-et-al-ricci-flow-part-iii:page-0522}
\uses{lem:chow-iii-i37-angle-decay, thm:chow-iii-g33-toponogov}
\dcref{appJ:I.38}
\group{Critical point theory}


Under the notation and hypotheses in the proof of Lemma I.37, if
\(r(x)\geq r_i/(1-\epsilon)\), \(d_i=d(x,x_i)\), and
\(c\in(0,\pi/2)\) is the angle bound occurring there, then
\[
 r(x)<d_i+r_i\sqrt{1-\epsilon^2\sin^2c}.
\]
\end{proposition}
\begin{proof}
\sketch
Consider the hinge at \(x_i\) formed by the minimizing geodesic from \(x_i\)
to \(x\) and the chosen minimizing geodesic from \(x_i\) to the basepoint.
Its angle \(\theta_i\) satisfies
\(\cos\theta_i>-\cos c\).  The hinge form of Toponogov comparison and the
Euclidean law of cosines imply
\[
 r(x)^2<d_i^2+r_i^2+2d_ir_i\cos c.
\]
Put \(a=\sin c\).  It is enough to prove
\[
 a^2\epsilon^2r_i
 \leq2d_i\left(\sqrt{1-a^2\epsilon^2}-\sqrt{1-a^2}\right),
\]
because adding this inequality to the preceding bound gives the square of
the desired estimate.  The triangle inequality hypothesis gives
\(r_i\leq(1-\epsilon)d_i/\epsilon\), so it suffices that
\[
 \frac{2(\sqrt{1-a^2\epsilon^2}-\sqrt{1-a^2})}
 {a^2\epsilon(1-\epsilon)}\geq1.
\]
After rationalizing the numerator, the left side becomes
\[
 \frac{2(1+\epsilon)}
 {\epsilon(\sqrt{1-a^2\epsilon^2}+\sqrt{1-a^2})},
\]
which is at least one because the denominator in parentheses is at most two
and \((1+\epsilon)/\epsilon\geq1\).  This proves the required estimate.
\end{proof}
